\chapter{Point Explosion}

\section{Supernova explosions}

Supernova explosions are very violent events which transfer a
significant amount of
energy to the interstellar medium (ISM) and are responsible for a large variety of physical
processes. We will not discuss the actual explosion mechanisms here which are thought to
be due to carbon deflagration of white dwarfs (type I) or the core collapse of massive stars
(type II) but we will follow the dynamical evolution of the supernova remnant (SNR), i.e.
the expanding cloud of hot gas in the ISM. This evolution can be divided into different
phases according to the dominant physical processes, and simplified models can be made
at least for the first stages, to get an idea of typical time scales, expansion velocities and
sizes.

\subsection{Free Expansion Phase}

The shock wave created by the explosion moves outwards into the (more or less homogeneous) interstellar gas at highly supersonic speed. During the first phase of the SNR
evolution the surrounding ISM has no influence on the expansion of the shock wave, i.e.
the pressure of the interstellar gas is negligible. Assuming that most of the supernova
energy $E_{\rm SN}$ (on the order of $10^{51}$ erg, excluding the part
carried away by neutrinos) is transformed into kinetic energy of the
ejected gas we can estimate the ejection velocity $u_e$ 

\beq
E_{\rm SN} = \frac{M_e u_e^2}{2} \quad \longrightarrow \quad u_e = \left(\frac{2E_{\rm SN}}{M_e}\right)^{1/2}
\eeq

\noindent where $M_e$ is the ejected mass. A supernova has kinetic
energy of about $10^{51}$ erg, and 1 $M_\sun$ of mass is ejected, so the
initial blast travels at $10^4$ km/s.


The shock radius scales as

\beq
R_s(t) = u_e t. 
\eeq

The schematic structure of the SNR at this phase can be described as follows: behind
the strong shock front which moves outwards into the ISM, compressed interstellar gas
is accumulating. This interstellar material is separated from the ejected stellar material
by a so-called contact discontinuity, i.e. a surface between two different materials with
similar pressure and velocity but different density (unlike a shock front where all three
quantities change). Behind the contact discontinuity the supernova
ejecta, traveling slower than the shock front, but faster than the
shocked interstellar gas, is racing up to shock against the dense
material upstream from the shock front. 

After some time the accumulated mass of the ISM compressed between the forward
shock and the contact discontinuity equals the ejected mass of stellar material, and it
will start to affect the expansion of the SNR. By definition, this is the end of the free
expansion phase, and the corresponding radius of the SNR, the
so-called sweep-up radius $R_{\rm SW}$, is defined by
\beq
M_e = \frac{4\pi}{3} R^3_{\rm SW} \rho_0 \quad \longrightarrow \quad
R_{\rm SW} = \left(\frac{3M_e}{4\pi \rho_0}\right)^{1/3}
\eeq

\noindent where $\rho_0$ is the initial density of the interstellar
medium. For $\rho_0=1.67\times 10^{-24}$ g/cm$^3$ (1 hydrogen atom per
cc), this yields $R_{\rm SW}$ of about 2 parsecs. This
radius is reached at the sweep-up time $t_{\rm SW} = R_{\rm
  SW}/u_e$. For $u_e=10^4$ km/s, this yields $t_{\rm SW}$ of roughly
250 years. 

Around the sweep-up time $t_{\rm SW}$ the structure of the SNR changes
due to the effects of the accumulated interstellar material. The second
the supernova ejecta has reached the shocked interstellar gas, leading
to a second shock front. While both shock fronts move forward, the
second shock is slower, so, from the reference frame of the ejecta it
is receeding; hence it is called {\it reverse} shock. The
reverse shock travel inwards through the supernova ejecta, heating
the ejected stellar gas to high temperatures. After some time a flat pressure structure
is established and the SNR enters the next phase where the expansion is driven by the
thermal pressure of the hot gas. Because the swept mass is now larger
than the ejecta mass, the free expansion is over and the shock starts
to decellerated. 

\subsection{Sedov-Taylor Phase}

After the passage of the reverse shock, the interior of the SNR is so hot that the energy
losses by radiation are very small (all atoms are ionized, no recombination). The following
pressure-driven expansion phase can therefore be regarded as adiabatic, the cooling of the
gas is only due to the expansion. Taylor (1950) and Sedov (1959) gave an exact self-
similar solution for such a pressure-driven explosion but we will restrict our analysis to
a simple estimate, that nevertheless captures the physics of the
problem.

\subsubsection{Dimensional analysis}

There’s an apocryphal story that Taylor managed to estimate the yield
of the first atomic bomb using declassified material. The story goes
that during the early days of nuclear testing, the precise yield of
the Trinity test was classified. However, some declassified footage of
the explosion had been made publicly available, showing the expanding
fireball in its early moments. Taylor then did the following
dimensional analysis.


Taylor reasoned that the problem was self-similar, as in
the evolution has to be of such nature as if some initial configuration
expands uniformly. So a subsequent configuration must appear like an
enlargement of a previous configuration. That is, the
explosion is homologous. 

Suppose an energy $E$ is suddenly released in an explosion producing
a spherical blast wave, which progresses in an ambient medium of
density $\rho$. We shall consider the pressure of the ambient medium
to be negligible compared to the pressure inside the blast wave (i.e.
$p \approx 0$). Let $\lambda$ be a scale parameter giving the size of the blast
wave at time $t$ after the explosion. Apart from being a monotonically
increasing function of time, the evolution of $\lambda$ may depend on $E$ and
$\rho$. Now the only way of combining $t$, $E$ and $\rho$ to get a dimension of
length is

\beq
\lambda = \left(\frac{Et^2}{\rho}\right)^{1/5}
\eeq


Let $r(t)$ be the radius of a shell of gas inside the spherical blast. For a
self-similar expansion, $r(t)$ has to evolve in the same way as $\lambda$. So we
can introduce a dimensionless distance parameter

\beq
\ksi = \frac{r}{\lambda} = r \left(\frac{\rho}{Et^2}\right)^{1/5}
\eeq

such that the value of this parameter for a particular shell of gas does
not change with time. Each shell of gas can therefore be labelled by a
particular value of $\ksi$. Let $\ksi_0$ correspond to the shock front so that the
radius of the spherical blast is given by

\beq
r_s(t) = \ksi_0 \left(\frac{Et^2}{\rho}\right)^{1/5}
\eeq

So the energy of the explosion is 

\beq
E = \frac{\rho}{t^2}\left(\frac{r_s}{\ksi_0}\right)^5 
\eeq

Assuming $\ksi_0=1$ (the actual factor is $\left(\sfrac{25}{4\pi}\right)^{1/5}\approx
1.15$, to be shown a posteriori), the energy of the explosion can be 
estimated from knowing $\rho$, and the radius $r_s$ of the blast a time $t$
after deflagration. As the story goes, Taylor saw the image shown in
\fig{fig:trinity} on the Time magazine, and estimated $r_s=75$ meters at
$t=0.006$ s. Knowing the air density $\rho=1.2\times 10^{-3}$\,g/cm$^3$.

\beq
E = \frac{1.2\times 10^{-3}}{\left(6\times
    10^{-3}\right)^2}\left(7.5\times 10^3\right)^5 \ {\rm erg}
\approx 8\times 10^{20} {\rm erg} 
\eeq

this is about 19 kilotons (1 kiloton $\approx 4.2\times 10^{19}$
erg), remarkably close to the actual yield of the Trinity, of 21
kilotons. As the story goes, Taylor published this result, while the
yield of the bomb was still classified information, generating ample embarrasment for the US
government. 

The story is probably apocryphal, as Taylor was part of the British
delegation to the Manhattan Project, and he was present at the Trinity
nuclear test. In any case, it illustrates how a simple analysis can
lead to accurate results. 

%contributing to the development
%of the bomb. atomic weapons. 
%Trinity Test and Taylor's Analysis:
%He was present at the Trinity nuclear test in 1945 and later, in 1950, published papers detailing his method of estimating the energy yield from photographs of the explosion. 

%While the details of this story have taken on a legendary quality, it remains a striking example of how pure physics can extract hidden truths from the world—even without direct access to secret information.

%Taylor was part of the British delegation to the Manhattan Project, contributing to the development of atomic weapons. 
%Trinity Test and Taylor's Analysis:
%He was present at the Trinity nuclear test in 1945 and later, in 1950, published papers detailing his method of estimating the energy yield from photographs of the explosion. 


%It was independently shown by Sedov (1946, 1959) and
%Taylor (1950a) that this problem can be solved if we assume that
%the blast wave evolves in a self-similar fashion. This means that
%the evolution has to be of such nature as if some initial configuration
%expands uniformly. So a subsequent configuration must appear like an
%enlargement of a previous configuration.

Continuing the dimensional analysis also yields the velocity of the
blast 

\beq
u_s(t) = \frac{dr_s}{dt}=\frac{2}{5}\frac{r_s}{t} = \frac{2r_s}{5t} =
\frac{2}{5}\ksi_0 \left(\frac{E}{\rho t^3}\right)^{1/2}
\eeq

so we find

\beqn
r_s &\propto& t^{2/5}\\
u_s &\propto& t^{-3/5}
\eeqn

A striking confirmation of the
$t^{2/5}$ dependence of size comes from the Trinity test. From a movie
of the explosion, Taylor (1950) plotted rhe graph shown in
\fig{fig:trinity-t25}, which shows
the variation of size with time, where the straight line corresponds to
the $t^{2/5}$ dependence. This plot provides a convincing proof that
self-similarity was a good assumption for this blast wave.

A typical supernova explosion ejects about $1M_\sun$ with an initial
velocity of about $10^4$ km/s, so that the energy input can be taken as
$E \approx 10^{51}$ erg. The blast wave expands in the interstellar medium with
typical density $\rho \approx 10^{-24}$ g\,cm$^3$. Again taking
$\ksi_0=1$, 

\beqn
r_s(t) &\approx& 2 \ {\rm pc} \ \left(\frac{t}{{\rm 100 \ yr}}\right)^{2/5} \\
u_s(t) &\approx& 5000 \ {\rm km/s} \ \left(\frac{t}{{\rm 100 \ yr}}\right)^{-3/5} 
\eeqn

This solution is valid after the free expansion phase, and before
radiative cooling starts to become important. 


\subsection{Ansatz solution}

In this section, we now find the appropriate value of $\ksi_0$. Taking
the gas pressure $p$ into account, the equation of motion for the expanding shell can be
formulated as force = pressure $\times$ area. Writing force as
$d(mv)/dt$, we can then write 

\beq
\frac{d}{dt}\left(\frac{4\pi}{3}R_s^3\rho_0 \dot{R}_s\right)=4\pi
R_s^2 p
\label{eq:radius-evolution}
\eeq

The pressure and internal energy $E$ of an ideal gas are related by

\beq
p  = \left(\gamma-1\right) \frac{E}{V}
\eeq

\noindent where $V$ denotes the volume. In the case of an adiabatic expansion (no losses of energy
due to radiation) we can set $E$ equal to the explosion energy $E_{\rm
  SN}$, and $V = 4\pi R^3_s /3$. Assuming $\gamma = 5/3$ we get

\beq
p  = \frac{E_{\rm SN}}{2\pi R_s^3}
\eeq

Making a power law ansatz for the shock radius

\beq
r_s = A t^\eta
\eeq

\noindent and inserting it into \eq{eq:radius-evolution} we obtain

\beq
\eta=\frac{2}{5} \quad {\rm and} \quad A  = \left(\frac{25E_{\rm SN}}{4\pi\rho_0}\right)^{1/5}
\eeq

\noindent i.e. the shock radius $r_s$ and the corresponding velocity $v_s$
behave like

\beqn
r_s(t) &=& \left(\frac{25E_{\rm SN}}{4\pi\rho_0}\right)^{1/5} t^{2/5}\\
u_s(t)&=&\frac{2}{5}\left(\frac{25E_{\rm SN}}{4\pi\rho_0}\right)^{1/5} t^{-3/5}
\eeqn

so we find

\beq
\ksi_0 = \left(\frac{25}{4\pi}\right)^{1/5} \approx 1.15
\eeq


\subsubsection{Temperature during Sedov-Taylor expansion}

The Sedov-Taylor phase of adiabatic expansion lasts until we can no longer ignore radiative
losses. %For that we need to compute the radiative loss 

%\beq
%\dot{E} = L = 4\pi r_s \sigma T^4
%\eeq

%and integrate it

%\beq
%E_{\rm loss} = \int \dot{E} dt 
%\eeq

%to find when $E_{\rm loss}$ is comparable to $E_{\rm SN}$ and no longer negligible.

The post-shock temperature of the gas can be derived using the Rankine-Hugoniot conditions (which express the basic conservation laws accross the shock front). In the limit
of a strong shock the post-shock pressure and density are given by

\beq
p_1 = \frac{2\rho_0 v_s^2}{\gamma+1} \quad {\rm and} \quad \frac{\rho_1}{\rho_0}=\frac{\gamma+1}{\gamma-1}
\eeq

Using these two conditions together with the ideal gas law

\beq
p_1 = \frac{k}{\mu m_u} \rho_1 T_1
\eeq

\noindent and assuming $\gamma = 5/3$ we obtain

\beq
T_1 = \frac{3\mu m_u}{16 k} u_s^2
\eeq

Inserting the expression for $u_s$ derived above we find
\beq
T_s(t) = T_1(t) = \frac{3\mu m_u}{100k}\left(\frac{25E_{\rm
      SN}}{4\pi\rho_0}\right)^{2/5}t^{-6/5}
\eeq

\noindent which describes the decrease of the post-shock temperature $T_s$ during the adiabatic,
pressure-driven expansion of the Sedov-Taylor phase.

\subsection{Cooling (snowplow) phase and dissipation into ISM}

As the SNR expands and cools adiabatically it will reach a critical temperature of about
$10^6$ K. At this temperature the ionized atoms start to capture free electrons and they can
lose their excitation energy by radiation. The radiative losses of energy become significant
putting an end to the adiabatic expansion of the SNR. Due to the efficient radiative cooling
the thermal pressure in the post-shock region decreases and the expansion slows down.
The SNR is entering the so-called snowplow phase (called snowplow because more and more interstellar gas
is accumulated until the swept-up mass is much larger than the ejected stellar material).

Finally the shell breaks up into individual clumps, probably due to a Rayleigh-Taylor
instability (hot thin gas, the ejecta, is pushing cool dense gas) and the SNR disperses into the ISM
as the expansion velocity decreases and becomes subsonic, with values typical of the interstellar gas.

The snowplow phase of a supernova remnant’s evolution is considered
momentum-conserving because, at this stage, the shock wave from the
supernova explosion has lost enough energy to radiation that the
expansion is primarily driven by the ram pressure (momentum) of the
ejected material rather than by thermal pressure. At this stage

\beq
mu={\rm const}
\eeq

Since $m \sim \rho r^3$ and $u=dr/dt$,

\beq
\rho r^3 \frac{dr}{dt} = {\rm const}
\eeq

solving for $r$ yields

\beq
r \propto t^{1/4}
\eeq

this phase lasts until the velocity drops to subsonic.


